% =============================================================================
% Block 4 --- Aula 4: consumption and saving (two-period problem, Euler, Ricardian equivalence, credit constraints).
% Weighted toward Prova 01 error cluster 3 (list every choice variable and every income source).
\bloco{4}{Consumption and saving}{Aula 4 \,\textbullet\, Kurlat ch.~6 \,\textbullet\, 10 points}

\questao{(4 points) \textit{Source: Kurlat, P. (2020), \emph{A Course in Modern Macroeconomics}, Exercise 6.4 ``Consumption and
Income'', pp.~122--123 --- statement reproduced verbatim, all items (a)--(d).}}

\noindent Suppose there are two households in the economy. Each of them solves a special case of Exercise 6.1, where
$\tau_1=\tau_2=a_0=0$, $\beta=1$ and $r=0$.
\begin{enumerate}[label=(\alph*),leftmargin=2.2em,topsep=2pt,itemsep=1pt]
\item Solve for $c_1$ as a function of $y_1$ and $y_2$.
\item Suppose the incomes of each of the households are:
\begin{center}\small
\begin{tabular}{@{}lcc@{}}
\toprule
 & $y_1$ & $y_2$ \\
\midrule
Household A & 2 & 4 \\
Household B & 6 & 4 \\
\bottomrule
\end{tabular}
\end{center}
Compute $c_1$ and $c_2$ for each of the households.
\item Suppose an economist is trying to decide what is a reasonable model for consumption behavior and only has data for period 1. Is
the data supportive of the Keynesian view of the consumption function? Explain.
\item Suppose the same economist now looks at data for period 2 in addition to data for period 1. Is the data supportive of the
Keynesian view of the consumption function? Explain.
\end{enumerate}
\noindent\textit{Our note, not in the book:} Kurlat's Exercise 6.1 (p.~121) is $\max_{c_1,a,c_2}u(c_1)+\beta u(c_2)$ s.t.\
$a=a_0+y_1-\tau_1-c_1$ and $c_2=y_2-\tau_2+(1+r)a$, with $u(c)=\frac{c^{1-\sigma}}{1-\sigma}$.\par
\noindent\textit{Added to the book's items (not in the book):} in (c) and (d), fit the line $c=a+by$ that a Keynesian economist would
estimate and say what it does and does not show.
\resolespaco
\begin{sol}
\begin{notacao}
\item $c_t$, $y_t$: consumption and income in period $t\in\{1,2\}$. $\beta=1$: discount factor. $r=0$: interest rate.
\item $u$: strictly increasing and strictly concave period utility (CRRA with curvature $\sigma>0$ in Exercise~6.1).
\item $a$, $b$: intercept and slope of the Keynesian line $c=a+by$; $\bar y$, $\bar c$: sample means.
\end{notacao}
\passo{(a)}{Exercise~6.1 with $\tau_1=\tau_2=a_0=0$ gives the Euler equation and the budget constraint:}
\begin{align}
u'(c_1)&=\beta(1+r)\,u'(c_2), \label{b4q10:euler}\\
c_1+\frac{c_2}{1+r}&=y_1+\frac{y_2}{1+r}. \label{b4q10:bc}
\end{align}
Substituting $\beta=1$ and $r=0$ into \eqref{b4q10:euler}:
\begin{align*}
u'(c_1)&=(1)(1+(0))\,u'(c_2)\\
u'(c_1)&=u'(c_2).
\end{align*}
$u$ is strictly concave, so $u'$ is strictly decreasing and
$c_1=c_2$ for any $\sigma$. Substituting $c_2=c_1$ into \eqref{b4q10:bc} with $r=0$:
\begin{align}
c_1+\frac{(c_1)}{1+(0)}&=y_1+\frac{y_2}{1+(0)}\nonumber\\
2c_1&=y_1+y_2\nonumber\\
c_1=c_2&=\frac{y_1+y_2}{2}. \label{b4q10:c}
\end{align}
\passo{(b)}{Using \eqref{b4q10:c} for each household:}
\begin{align*}
\text{A:}\quad c_1=c_2&=(2+4)/2=3,\qquad y_1-c_1=2-3=-1,\\
\text{B:}\quad c_1=c_2&=(6+4)/2=5,\qquad y_1-c_1=6-5=1 .
\end{align*}
A borrows $1$ in period 1; B saves $1$.
\passo{(c) period-1 data}{The line through $(y,c)=(2,3)$ and $(6,5)$:}
\begin{align*}
b&=\frac{5-3}{6-2}=0.5,\\
a&=3-0.5\times2=2 .
\end{align*}
So $c=2+0.5\,y$: a positive intercept, a marginal propensity of $0.5<1$, and an average propensity that falls with income
($3/2=1.5$ against $5/6=0.83$). On its face this is Keynes's consumption function. But the slope is not a behavioural response:
income differences in period 1 are partly temporary, and by \eqref{b4q10:c} each household consumes its \emph{average} income.
\passo{(d) both periods}{In period 2 both incomes are $4$ while consumption is $3$ and $5$: current income does not explain
consumption at all. Over time A's income doubles ($2\to4$) and B's falls ($6\to4$), and neither household changes $c$: the
time-series propensity to consume out of these anticipated income changes is zero. Pooling the four observations, the least-squares
line is}
\begin{equation}
b=\frac{\mathrm{cov}(y,c)}{\mathrm{var}(y)},\qquad a=\bar c-b\,\bar y. \label{b4q10:ols}
\end{equation}
With $y=(2,6,4,4)$ and $c=(3,5,3,5)$: $\bar y=4$, $\bar c=4$, $\mathrm{cov}(y,c)=1$, $\mathrm{var}(y)=2$. Using \eqref{b4q10:ols}:
\begin{align*}
b&=1/2=0.5,\\
a&=4-0.5\times4=2 .
\end{align*}
The pooled line is still $c=2+0.5\,y$, yet it fits no household's behaviour: the data reject the Keynesian function and support the
permanent-income view of \eqref{b4q10:c}, where $c$ depends on lifetime resources ($y_1+y_2=6$ for A and $10$ for B).
\intuicao{A cross-section confounds permanent and temporary income; consumption responds to the first, so a Keynesian slope fitted
on one year of data measures how much income variation is permanent, not how households react to income.}
\refbib{Kurlat (2020), \S6.1 (p.~103) and \S6.2 (p.~105), eqs.~(6.2.10)--(6.2.11). Consistent with the project's worked solution
(Problem 6.4, folder \texttt{Resolucao}); the book prints no answer.}
\rubrica{
\item (a) 1; (b) 1; (c) 1; (d) 1. Calling the period-1 slope the households' MPC without qualification: max.\ 0.5 pt in (c).
}
\end{sol}

% -----------------------------------------------------------------------------
\questao{(5 points) Discursive. \emph{Choice variables $\times$ Ricardian equivalence $\times$ credit constraints.} A Brazilian worker lives two
periods (working life and retirement, about 30 years each) with utility $\ln c_1+\beta\ln c_2$, $\beta=0.6$. In period 1 she holds initial
wealth $a_0=10$, earns labour income $w_1=40$ (hours are fixed), receives dividends $\Pi_1=10$ from the family firm she owns and pays a
lump-sum tax $\tau_1=10$. In period 2 she receives income $y_2=90$ and pays $\tau_2=10$. She borrows or lends $a$ at the gross rate
$1+r=2$ (about $2.3\%$ a year over 30 years).}

\textbf{a)} (2 points) List every choice variable and every source of income. Write the two period budget constraints and the
intertemporal one, set up the Lagrangian, derive the first-order conditions and the Euler equation, and find $c_1$, $c_2$ and $a$.
\resolespaco
\begin{sol}
\begin{notacao}
\item $c_1$, $c_2>0$: consumption in working life and in retirement. $a\in\mathbb R$: assets carried into period 2 ($a<0$ is borrowing).
\item $a_0=10$: initial wealth. $w_1=40$: labour income. $\Pi_1=10$: dividends. $y_2=90$: period-2 income. $\tau_1$, $\tau_2$: lump-sum taxes.
\item $\beta\in(0,1)$: discount factor, $0.6$. $r$: interest rate over the period, $1+r=2$.
\item $W$: lifetime resources in present value. $\lambda>0$: multiplier on the intertemporal budget.
\end{notacao}
\passo{choices and income}{Choice variables: $c_1$, $c_2$, $a$ (a negative $a$ is borrowing). Income: initial wealth $a_0$, labour income
$w_1$, profits $\Pi_1$ (she owns the firm), period-2 income $y_2$; taxes $\tau_1$, $\tau_2$ are paid out of them.}
\passo{constraints}{Resources equal uses in each period:}
\begin{align}
c_1+a&=a_0+w_1+\Pi_1-\tau_1, \label{b4q11:bc1}\\
c_2&=y_2-\tau_2+(1+r)\,a. \label{b4q11:bc2}
\end{align}
Solve \eqref{b4q11:bc2} for $a$:
\[
a=\frac{c_2-y_2+\tau_2}{1+r}.
\]
Substitute this into \eqref{b4q11:bc1}:
\[
c_1+\Big(\frac{c_2-y_2+\tau_2}{1+r}\Big)=a_0+w_1+\Pi_1-\tau_1 .
\]
Collect the terms in $c_2$ on the left:
\begin{equation}
c_1+\frac{c_2}{1+r}=a_0+w_1+\Pi_1-\tau_1+\frac{y_2-\tau_2}{1+r}\equiv W. \label{b4q11:ibc}
\end{equation}
With the data:
\[
W=10+40+10-10+\frac{90-10}{2}=50+40=90 .
\]
\passo{Lagrangian and FOCs}{Maximise utility subject to \eqref{b4q11:ibc}:}
\begin{equation}
\mathcal L=\ln c_1+\beta\ln c_2+\lambda\Big[W-c_1-\frac{c_2}{1+r}\Big]. \label{b4q11:L}
\end{equation}
Differentiate \eqref{b4q11:L} with respect to $c_1$ and $c_2$:
\begin{align}
\frac1{c_1}&=\lambda, \label{b4q11:foc1}\\
\frac{\beta}{c_2}&=\frac{\lambda}{1+r}. \label{b4q11:foc2}
\end{align}
Substitute $\lambda$ from \eqref{b4q11:foc1} into \eqref{b4q11:foc2} and rearrange:
\begin{align}
\frac{\beta}{c_2}&=\frac{(1/c_1)}{1+r}\nonumber\\
\frac{\beta}{c_2}&=\frac{1}{(1+r)\,c_1}\nonumber\\
c_2&=\beta(1+r)\,c_1. \label{b4q11:euler}
\end{align}
With the data, $c_2=0.6\times2\times c_1=1.2\,c_1$.
\passo{solution}{Substitute \eqref{b4q11:euler} into \eqref{b4q11:ibc}:}
\begin{align*}
c_1+\frac{\big(\beta(1+r)\,c_1\big)}{1+r}&=W\\
c_1+\beta c_1&=W\\
(1+\beta)\,c_1&=W,
\end{align*}
so
\begin{equation}
c_1=\frac{W}{1+\beta}=\frac{90}{1.6}=56.25. \label{b4q11:c1}
\end{equation}
Then \eqref{b4q11:euler} and \eqref{b4q11:bc1} give
\begin{align*}
c_2&=1.2\times(56.25)=67.5,\\
a&=10+40+10-10-(56.25)\\
&=50-56.25=-6.25 .
\end{align*}
Check with \eqref{b4q11:bc2}: $c_2=80+2\times(-6.25)=67.5$.
\intuicao{She consumes the fraction $1/(1+\beta)$ of lifetime wealth whatever its timing; her income is back-loaded ($y_2$ is large), so
smoothing makes her a borrower.}
\refbib{Kurlat (2020), \S6.2 (p.~105) and Exercise~6.1 (p.~121). The closed form of $c_1$ is the $\sigma=1$ case of the instructor's
solution to Lista~3, Q1(a) (\texttt{Listas/solu\c{c}\~oes do instrutor/PSET 3 (1).pdf}).}
\rubrica{
\item Choice variables and income sources 0.5; constraints and $W$ 0.5; FOCs and Euler 0.5; $c_1,c_2,a$ with check 0.5.
Profits $\Pi_1$ or $a_0$ omitted from the budget (the Prova~01 Q2(c) error): max.\ 1 pt. Discounting $c_2$ by $(1+r)$ in the wrong direction: max.\ 0.5 pt.
}
\end{sol}

\textbf{b)} (1.5 points) The government cuts $\tau_1$ to $0$, finances the gap by selling bonds and repays them with interest by raising $\tau_2$ to
$30$. Find the new $c_1$, $c_2$ and $a$, the change in private, public and national saving, and the assumptions the result needs.
\resolespaco
\begin{sol}
\begin{notacao}
\item Symbols as in (a); a prime marks values after the tax change ($\tau_1'=0$, $\tau_2'=30$).
\item $S^p$, $S^g$, $S$: private, public and national saving in period 1, with $S=S^p+S^g$.
\end{notacao}
\passo{wealth}{Using \eqref{b4q11:ibc} with the new taxes:}
\[
W'=10+40+10-0+\frac{90-30}{2}=60+30=90=W .
\]
The tax cut of $10$ is matched by a tax rise of $20=10\times(1+r)$, the same in present value.
\passo{choices}{$W$, $\beta$ and $r$ are unchanged, so \eqref{b4q11:c1} and \eqref{b4q11:euler} give $c_1=56.25$ and $c_2=67.5$. By \eqref{b4q11:bc1}:}
\begin{align*}
a'&=10+40+10-0-(56.25)\\
&=60-56.25=3.75 ,
\end{align*}
so $a$ rises by exactly $10$, the bonds she buys.
\passo{saving}{The changes in period 1:}
\begin{align*}
\Delta S^p&=\Delta a=3.75-(-6.25)=10,\\
\Delta S^g&=\Delta\tau_1=-10\ \text{(the deficit)},\\
\Delta S&=10-10=0 .
\end{align*}
National saving is unchanged: Ricardian equivalence.
\passo{assumptions}{Lump-sum taxes; the government and the household borrow at the same $r$; the taxpayer who gets the cut pays the
later tax (same horizon); she can borrow and lend freely at $r$.}
\refbib{Kurlat (2020), \S6.2, ``Taxes and Ricardian Equivalence''; Exercise~6.5 (p.~123).}
\rubrica{\item Wealth 0.5; choices 0.5; saving and assumptions 0.5. Claiming $c_1$ rises: max.\ 0.5 pt.}
\end{sol}

\textbf{c)} (1.5 points) Now she cannot borrow: $a\ge0$. Solve her problem before and after the tax change of (b). Does Ricardian equivalence hold?
What is her marginal propensity to consume out of the tax cut?
\resolespaco
\begin{sol}
\begin{notacao}
\item Symbols as in (a). $\mu\ge0$: multiplier on the borrowing constraint $a\ge0$.
\end{notacao}
\passo{the problem}{Substitute \eqref{b4q11:bc1} and \eqref{b4q11:bc2} for $c_1$ and $c_2$, so that $a$ is the only choice:}
\begin{equation}
\mathcal L=\ln\big(a_0+w_1+\Pi_1-\tau_1-a\big)+\beta\ln\big(y_2-\tau_2+(1+r)a\big)+\mu\,a. \label{b4q11:Lc}
\end{equation}
Differentiate \eqref{b4q11:Lc} with respect to $a$:
\begin{equation}
-\frac1{c_1}+\frac{\beta(1+r)}{c_2}+\mu=0,\qquad \mu\ge0,\qquad \mu\,a=0. \label{b4q11:foca}
\end{equation}
\passo{before}{The unconstrained choice $a=-6.25$ violates $a\ge0$, so the constraint binds: $a=0$, and \eqref{b4q11:bc1}--\eqref{b4q11:bc2} give}
\begin{align*}
c_1&=a_0+w_1+\Pi_1-\tau_1-(0)\\
&=10+40+10-10=50,\\
c_2&=y_2-\tau_2+(1+r)(0)\\
&=90-10=80 .
\end{align*}
Solve \eqref{b4q11:foca} for $\mu$:
\begin{align*}
\mu&=\frac1{c_1}-\frac{\beta(1+r)}{c_2}\\
&=\frac1{50}-\frac{1.2}{80}=0.02-0.015=0.005>0 .
\end{align*}
The Euler equation holds as an inequality, $u'(c_1)>\beta(1+r)u'(c_2)$.
\passo{after}{Period-1 resources are $60$ and the unconstrained choice of (b) gives $a=3.75\ge0$: the constraint no longer binds
($\mu=0$ in \eqref{b4q11:foca}, which is then the Euler equation \eqref{b4q11:euler}), so $c_1=56.25$, $c_2=67.5$.}
\passo{equivalence fails}{$c_1$ rises from $50$ to $56.25$ out of a tax cut that leaves $W$ unchanged:}
\[
\text{MPC}=\frac{56.25-50}{10}=0.625 .
\]
The cut works as a loan from the government that the market would not give her.
\intuicao{Credit-constrained households consume current income, so the timing of taxes matters for them and a debt-financed transfer
raises demand.}
\refbib{Kurlat (2020), Exercise~6.5 \emph{Credit Constraints and Ricardian Equivalence} (p.~123); \S6.2, ``Taxes and Ricardian Equivalence''.}
\rubrica{
\item Before 0.5 (with the inequality); after 0.5; MPC and conclusion 0.5. Imposing the Euler equality while the constraint binds: max.\ 0.5 pt.
}
\end{sol}

% -----------------------------------------------------------------------------
\questao{(1 point) Multiple choice. A household with $u(c)=c^{1-\sigma}/(1-\sigma)$, $\sigma>0$, has income only in period 1 ($y_1>0$, $y_2=0$, no
taxes, no initial wealth) and saves at rate $r$. When $r$ rises, $\partial c_1/\partial r$:}
\begin{alternativas}
\item is negative for every $\sigma>0$: a higher $r$ makes future consumption cheaper, so she saves more.
\item is positive iff $\sigma>1$: for a saver, the income effect of $r$ then outweighs the substitution effect.
\item is positive iff $\sigma<1$: an elasticity of intertemporal substitution above one makes the income effect dominate.
\item is zero for every $\sigma>0$: with $y_2=0$ there is no future income for a change in $r$ to revalue.
\item is negative iff $\sigma>1$: $\sigma$ is the elasticity of intertemporal substitution, so a large $\sigma$ makes substitution dominate.
\end{alternativas}
\gabarito{B}
\begin{sol}
\begin{notacao}
\item $c_1$, $c_2$: consumption; $y_1>0$: period-1 income; $\beta\in(0,1)$: discount factor; $r>-1$: interest rate;
$\sigma>0$: curvature of $u$, so $1/\sigma$ is the EIS; $D$: defined in \eqref{b4q12:c1}.
\end{notacao}
\passo{closed form}{The CRRA Euler equation and the budget constraint with $y_2=0$:}
\begin{align}
c_2&=[\beta(1+r)]^{1/\sigma}c_1, \label{b4q12:euler}\\
c_1+\frac{c_2}{1+r}&=y_1. \label{b4q12:bc}
\end{align}
Substitute \eqref{b4q12:euler} into \eqref{b4q12:bc}:
\begin{align}
c_1+\frac{\big([\beta(1+r)]^{1/\sigma}c_1\big)}{1+r}&=y_1\nonumber\\
c_1\big[1+\beta^{1/\sigma}(1+r)^{\frac1\sigma-1}\big]&=y_1\nonumber\\
c_1=\frac{y_1}{D},\qquad D&\equiv1+\beta^{1/\sigma}(1+r)^{\frac1\sigma-1}. \label{b4q12:c1}
\end{align}
\passo{sign}{Differentiate \eqref{b4q12:c1} with respect to $r$:}
\begin{align*}
\frac{\partial c_1}{\partial r}&=-\frac{y_1}{D^2}\,\frac{\partial D}{\partial r}\\
&=-\frac{y_1}{D^2}\,\Big(\beta^{1/\sigma}\Big(\frac1\sigma-1\Big)(1+r)^{\frac1\sigma-2}\Big)\\
&=\frac{y_1\beta^{1/\sigma}(1+r)^{\frac1\sigma-2}}{D^2}\cdot\frac{\sigma-1}{\sigma}.
\end{align*}
The first factor is positive, so the derivative is
positive iff $\sigma>1$, zero at $\sigma=1$ (log), negative iff $\sigma<1$. The EIS is $1/\sigma$; with $\sigma>1$ it is below one and the income
effect wins. A is the substitution effect alone; C reverses the condition; D ignores that $r$ also raises her interest income; E reverses
the sign and calls $\sigma$ the EIS (two mistakes).
\refbib{Kurlat (2020), \S6.2 (p.~105) and Exercises~6.1(e), 6.3 (pp.~121--122). Confirmed by the instructor's solution to Lista~3,
Q1(c): $\partial c_1/\partial r$ has the sign of $(\sigma-1)/\sigma$ (\texttt{PSET 3 (1).pdf}).}
\end{sol}
